\documentclass[a4paper,12pt]{article}

\usepackage[french]{babel}
\usepackage[utf8]{inputenc}
\usepackage{graphicx}
\usepackage{url}

\title{TD 7 : Mini projets : mise en situation}
\author{EL FHAL Dalia \and AJOUIRJA Hajar \and MIFTARI Eris}
\date{Mars 2026}

        




\maketitle



\resume{
Ce travail est une réalisation pour notre cours de gestion de projet de notre deuxième année de Licence.C'est une mise en situation des outils de gestion de projet à travers des cas pratiques. 

   
}



\begin{document}



\section{Introduction}  
Notre groupe de travail est composé de 3 étudiants : Hajar Ajouirja, Dalia El Fhal et Eris Miftari. Nous faisons un porjet Latex.

Notre projet est sur GitHub sous le nom de : Projet-Mise-En-Situation

Pour la répartition des taches ça a été assez complexe puisque certains membres ont rencontré des soucis donc on a dû adapter aux besoins.




\section{Description des mises en situations}
\subsection{nom de la partie}



\subsection{autre nom}












\section{Conclusion}

Ce mini-projet a constitué une expérience enrichissante pour notre groupe, tant sur le plan méthodologique que sur le plan humain. Il nous a permis de mieux comprendre l’importance de l’organisation, de la communication et de la répartition des tâches dans la conduite d’un projet. Nous avons pu nous familiariser davantage avec l’utilisation de l’outil LaTeX ainsi qu’avec le travail collaboratif à distance,via GitHub.

Cependant, nous regrettons de ne pas avoir anticipé certaines difficultés techniques, en particulier le fait que le projet ne fonctionnait pas toujours correctement sur nos ordinateurs personnels (MacOS). De plus, nous avons commencé le travail relativement tard dans la semaine n'ayant eu le TD que jeudi après-midi, ce qui a réduit notre temps de production. Nos indisponibilités dû à nos jobs étudiants en fin de semaine ont également limité les possibilités d’échanges et d’amélioration du rendu final.

Malgré ces contraintes, ce projet nous a permis de tirer des enseignements importants pour nos futurs travaux universitaires. Il nous a notamment fait prendre conscience de la nécessité d’anticiper les imprévus, de mieux planifier notre temps et de débuter plus tôt afin de produire un travail plus approfondi et de meilleure qualité.




\end{document}
